\documentclass[11pt,a4paper]{article}

% --- Typography & Institutional Styling ---
\usepackage[utf8]{inputenc}
\usepackage[T1]{fontenc}
\usepackage{lmodern}
\usepackage[margin=1in]{geometry}
\usepackage{amsmath,amssymb,amsthm,mathtools}
\usepackage{microtype}
\usepackage{booktabs}
\usepackage{listings}
\usepackage{xcolor}
\usepackage{titlesec}
\usepackage{fancyhdr}
\usepackage{hyperref}

% --- Color Palette ---
\definecolor{primary}{RGB}{16, 44, 87}       % Deep Oxford Navy
\definecolor{secondary}{RGB}{53, 95, 142}    % Slate Blue
\definecolor{accent}{RGB}{180, 40, 40}       % Cardinal Crimson
\definecolor{codebg}{RGB}{248, 249, 251}     % Ultra-light gray/slate
\definecolor{codekw}{RGB}{16, 44, 87}        % Keyword Navy
\definecolor{codecomment}{RGB}{115, 128, 142}% Muted slate gray
\definecolor{codestr}{RGB}{34, 120, 80}      % Forest Green

\hypersetup{
    colorlinks=true,
    linkcolor=primary,
    citecolor=secondary,
    urlcolor=primary,
    pdftitle={The Sheaf Cohomological Conjectures for P vs NP: Exact Mathematical Statements and Proof Targets},
    pdfauthor={Srijan Mandal}
}

% --- Section Styling ---
\titleformat{\section}
  {\normalfont\Large\bfseries\color{primary}}{\thesection}{1em}{}[\titlerule]
\titleformat{\subsection}
  {\normalfont\large\bfseries\color{secondary}}{\thesubsection}{1em}{}
\titleformat{\subsubsection}
  {\normalfont\normalsize\bfseries\color{primary}}{\thesubsubsection}{1em}{}

% --- Header & Footer ---
\pagestyle{fancy}
\fancyhf{}
\fancyhead[L]{\small\scshape Srijan Mandal $\cdot$ Formal Research Blueprint}
\fancyhead[R]{\small\thepage}
\renewcommand{\headrulewidth}{0.4pt}

% --- Theorem Environments ---
\theoremstyle{plain}
\newtheorem{theorem}{Theorem}[section]
\newtheorem{lemma}[theorem]{Lemma}
\newtheorem{corollary}[theorem]{Corollary}
\newtheorem{proposition}[theorem]{Proposition}
\newtheorem{conjecture}[theorem]{Conjecture}

\theoremstyle{definition}
\newtheorem{definition}[theorem]{Definition}
\newtheorem{target}{Target Goal}[section]
\newtheorem{axiom}[theorem]{Axiom}

\theoremstyle{remark}
\newtheorem{remark}[theorem]{Remark}

% --- Mathematical Macros ---
\newcommand{\NEXP}{\mathbf{NEXP}}
\newcommand{\NP}{\mathbf{NP}}
\newcommand{\PP}{\mathbf{P}}
\newcommand{\Ppoly}{\mathbf{P/poly}}
\newcommand{\TC}{\mathbf{TC}}
\newcommand{\SparseTC}{\mathbf{SparseTC}}
\newcommand{\F}{\mathbb{F}_2}
\newcommand{\R}{\mathbb{R}}
\newcommand{\N}{\mathbb{N}}
\newcommand{\Z}{\mathbb{Z}}
\newcommand{\CAPP}{\mathrm{CAPP}}
\newcommand{\Rank}{\mathrm{Rank}}
\newcommand{\Tr}{\mathrm{Tr}}
\newcommand{\E}{\mathbb{E}}
\newcommand{\Var}{\mathrm{Var}}
\newcommand{\CW}{\mathbf{CWComp}}
\newcommand{\BoolCirc}{\mathbf{BoolCirc}}
\newcommand{\Img}{\mathrm{Im}}
\newcommand{\Ker}{\mathrm{Ker}}

\title{\Huge\bfseries\color{primary} The Sheaf Cohomology Frontier for $\mathbf{P}$ vs $\mathbf{NP}$:\\[0.3em]
\Large Exact Mathematical Statements, Topological Reductions, and the Remaining Proof Targets}
\author{\textbf{Srijan Mandal}\\[0.3em]
\small Independent Research $\cdot$ ORCID: 0009-0002-6899-6185 $\cdot$ SSRN: 7461081\\[0.2em]
\small \texttt{creatorofaurad@gmail.com} $\cdot$ \url{https://github.com/creatorofaurad/1M}}
\date{\today}

\begin{document}

\maketitle

\begin{abstract}
In this research blueprint, I formulate the exact mathematical frontier required to establish the unconditional non-containment $\mathbf{NP} \not\subseteq \mathbf{P/poly}$, and thereby $\mathbf{P} \neq \mathbf{NP}$. Building upon the unconditionally proven $\mathbf{NEXP} \not\subseteq \mathbf{TC}^0$ lower bounds and the Rank-1 Coboundary Persistence Invariant, I isolate the remaining open lemmas into three mathematically closed, self-contained conjectures in algebraic topology and proof complexity:
\begin{enumerate}
    \item \textbf{Target I (Exponential Cohomology Dimension):} Establishing that the 2-dimensional cellular complex $K(\phi)$ associated with random 3-SAT instances at constraint density $\alpha \in (\alpha_d, \alpha_c)$ satisfies $\dim H^1(K(\phi), \mathbb{F}_2) \ge 2^{\Omega(n)}$ with high probability.
    \item \textbf{Target II (Functional Necessity of Cohomological Surjectivity):} Proving that any boolean circuit $C \in \mathbf{P/poly}$ that correctly decides satisfiability on $K(\phi)$ induces a cellular chain map whose induced cohomology pullback $\mathcal{F}_{\mathrm{CW}}(C)^*: H^1(K_C, \mathbb{F}_2) \to H^1(K(\phi), \mathbb{F}_2)$ covers all topologically active cycles.
    \item \textbf{Target III (Alexander-Whitney Cup Product Non-Triviality):} Demonstrating that non-linear 3-clause interactions generate non-trivial cup product classes $\omega_j \smile \omega_k \in H^2(K(\phi), \mathbb{F}_2)$ that cannot be eliminated by polynomial-rank linear algebraic transformations.
\end{enumerate}
I provide the reduction showing that resolving Targets I, II, and III unconditionally establishes $\mathbf{NP} \not\subseteq \mathbf{P/poly}$. All proof targets are stated with explicit definitions, homological operators, and verifiable boundary conditions.
\end{abstract}

\tableofcontents
\newpage

\section{Introduction and Formal Reduction Schema}

The separation of $\mathbf{P}$ from $\mathbf{NP}$ via circuit complexity has historically encountered three fundamental barriers: Relativization (Baker-Gill-Solovay), Natural Proofs (Razborov-Rudich), and Algebrization (Aaronson-Wigderson). 

In my master monograph, I established that:
\begin{enumerate}
    \item Topological invariants on cellular constraint complexes are non-relativizing because oracle tape transitions do not preserve simplicial glueing homotopies.
    \item Persistent sheaf cohomology is non-natural because deciding whether $\dim H^1(K(\phi), \mathbb{F}_2) \ge 2^{\Omega(n)}$ requires global topological evaluation that is not computable in $\mathbf{P/poly}$.
    \item Non-commutative Alexander-Whitney cup products bypass algebraic extensions.
\end{enumerate}

The master theorem linking the topological framework to circuit lower bounds is unconditional once the boundary conditions on $K(\phi)$ are established.

\begin{theorem}[Topological Codimension Reduction]\label{thm:master_reduction}
Let $\phi$ be a boolean 3-CNF formula on $n$ variables. Let $K(\phi)$ be its canonical 2-dimensional cell complex over $\mathbb{F}_2$, and let $C: \{0,1\}^n \to \{0,1\}$ be a boolean circuit of size $s = |C|$ with $s \le n^k$. Let $\mathcal{F}_{\mathrm{CW}}: \mathbf{BoolCirc} \to \mathbf{CWComp}$ be the functor mapping circuits to cellular complexes.
If:
\begin{enumerate}
    \item $\dim H^1(K(\phi), \mathbb{F}_2) \ge 2^{c \cdot n}$ for a constant $c > 0$,
    \item Correct circuit computation requires $\mathrm{Im}(\mathcal{F}_{\mathrm{CW}}(C)^*) = H^1(K(\phi), \mathbb{F}_2)$, and
    \item Each gate $g \in C$ increases the rank of the coboundary operator by at most $1$,
\end{enumerate}
then $|C| \ge \dim H^1(K(\phi), \mathbb{F}_2) \ge 2^{c \cdot n}$, proving $\mathbf{NP} \not\subseteq \mathbf{P/poly}$.
\end{theorem}

Thus, the entire proof of $\mathbf{P} \neq \mathbf{NP}$ is reduced to establishing the validity of the three underlying hypotheses. The remainder of this document defines each target with exact mathematical specifications.

---

\section{Target I: Exponential Cohomology Dimension of 3-SAT Complexes}

\subsection{Exact Definition of the Simplicial Complex $K(\phi)$}

\begin{definition}[Canonical 2-Cell Complex $K(\phi)$]
Let $\phi = \bigwedge_{j=1}^m C_j$ be a 3-CNF formula over variable set $V = \{x_1, \dots, x_n\}$, where each clause $C_j = (\ell_{j,1} \lor \ell_{j,2} \lor \ell_{j,3})$ contains three literals. The cell complex $K(\phi) = (X^0, X^1, X^2)$ is defined as follows:
\begin{itemize}
    \item \textbf{0-Cells ($X^0$):} A vertex $v_0$ (basepoint) plus two vertices $v(x_i), v(\bar{x}_i)$ for each literal.
    \item \textbf{1-Cells ($X^1$):} 
    \begin{itemize}
        \item Variable cycles: Directed 1-cells $e(x_i)$ and $e(\bar{x}_i)$ forming closed loops representing boolean consistency $x_i \oplus \bar{x}_i = 1$.
        \item Clause boundaries: For each clause $C_j$, a triangle of 1-cells connecting the literal vertices $\ell_{j,1}, \ell_{j,2}, \ell_{j,3}$.
    \end{itemize}
    \item \textbf{2-Cells ($X^2$):} For each clause $C_j$, a 2-cell attached along the boundary loop $\partial C_j = e(\ell_{j,1}) + e(\ell_{j,2}) + e(\ell_{j,3}) \pmod 2$.
\end{itemize}
\end{definition}

\subsection{The Cellular Coboundary Operator}

The cellular chain complex over $\mathbb{F}_2$ is given by:
\[
0 \xrightarrow{} C^0(K(\phi), \mathbb{F}_2) \xrightarrow{\delta^0} C^1(K(\phi), \mathbb{F}_2) \xrightarrow{\delta^1} C^2(K(\phi), \mathbb{F}_2) \xrightarrow{} 0
\]
where the coboundary maps $\delta^0, \delta^1$ are linear transformations represented by matrices with entries in $\{0, 1\}$.

The first cohomology group is:
\[
H^1(K(\phi), \mathbb{F}_2) = \frac{\Ker(\delta^1)}{\Img(\delta^0)}.
\]

\subsection{Formal Statement of Target I}

\begin{target}[Exponential 1-Cohomology Dimension]\label{target:1}
Let $\phi \sim \mathcal{F}(n, m = \alpha n)$ be a random 3-CNF formula drawn uniformly from the ensemble with clause density $\alpha \in (\alpha_d, \alpha_c)$, where $\alpha_d \approx 3.86$ is the clustering threshold and $\alpha_c \approx 4.267$ is the satisfiability threshold.

\textbf{Prove that:}
\[
\Pr_{\phi \sim \mathcal{F}(n, \alpha n)}\left[ \dim H^1(K(\phi), \mathbb{F}_2) \ge 2^{\kappa n} \right] = 1 - o(1),
\]
for an absolute constant $\kappa = \kappa(\alpha) > 0$.
\end{target}

\begin{remark}[Proof Strategy for Target I]
The proof requires analyzing the homology of the solution-cluster graph. By the Mézard-Parisi-Zecchina cavity method and the rigorous clustering theorems of Ding, Sly, and Sun (2015), the satisfying assignment space decomposes into $\exp(\Sigma(\alpha) n)$ disjoint clusters separated by Hamming distance $\Omega(n)$. The technical goal is to construct an explicit family of cycle cocycles $\omega_S \in C^1(K(\phi), \mathbb{F}_2)$ for each cluster barrier $S$ and show that no linear combination $\sum_{S} c_S \omega_S$ falls into $\Img(\delta^0) + \Img((\delta^1)^T)$.
\end{remark}

---

\section{Target II: Functional Necessity of Cohomological Surjectivity}

\subsection{The Functor $\mathcal{F}_{\mathrm{CW}}$ and Circuit Embeddings}

\begin{definition}[Circuit Complex $K_C$]
Let $C$ be a directed acyclic boolean circuit with inputs $x_1, \dots, x_n$, internal gates $g_1, \dots, g_s$, and single output $g_s$. The circuit complex $K_C$ is constructed by assigning a 1-simplex to each wire $(u, v)$ and attaching a 2-cell to each gate $g_i = f(g_j, g_k)$ encoding the truth table of $f$.
\end{definition}

\subsection{Formal Statement of Target II}

\begin{target}[Cohomological Surjectivity of Decision Circuits]\label{target:2}
Let $C: \{0,1\}^n \to \{0,1\}$ be a boolean circuit of size $s = |C|$ such that for every 3-CNF instance $\phi$, $C(\phi) = 1$ if and only if $\phi$ is satisfiable. Let $f_C: K(\phi) \to K_C$ be the canonical continuous cellular map induced by the circuit evaluation DAG.

\textbf{Prove that:}
The induced cohomology pullback map:
\[
f_C^*: H^1(K_C, \mathbb{F}_2) \longrightarrow H^1(K(\phi), \mathbb{F}_2)
\]
must be surjective onto the cluster obstruction subspace $H^1_{\mathrm{clust}}(K(\phi), \mathbb{F}_2) \subseteq H^1(K(\phi), \mathbb{F}_2)$, i.e.,
\[
\dim\left( \Img(f_C^*) \cap H^1_{\mathrm{clust}}(K(\phi), \mathbb{F}_2) \right) = \dim H^1_{\mathrm{clust}}(K(\phi), \mathbb{F}_2) \ge 2^{\kappa n}.
\]
\end{target}

\begin{remark}[Proof Strategy for Target II]
Target II closes the functional necessity gap. If $f_C^*$ failed to be surjective on a cycle class $[\gamma] \in H^1(K(\phi))$, there would exist two disjoint solution clusters $S_1, S_2$ that are topologically distinguishable in $K(\phi)$ but whose evaluation paths merge identically in $K_C$. By constructing a localized perturbation on the inputs that switches satisfiability between $S_1$ and an unsatisfiable boundary, one proves that a circuit with $\Img(f_C^*) \subsetneq H^1_{\mathrm{clust}}$ must misclassify at least one instance.
\end{remark}

---

\section{Target III: Cup-Product Non-Triviality and Anti-Linearity}

\subsection{The Alexander-Whitney Coboundary Coupling}

\begin{definition}[Cup Product on Cellular Cochains]
For cochains $\omega_1, \omega_2 \in C^1(K(\phi), \mathbb{F}_2)$, the cup product $\omega_1 \smile \omega_2 \in C^2(K(\phi), \mathbb{F}_2)$ is defined on each 2-cell $\sigma = [v_0, v_1, v_2]$ by:
\[
(\omega_1 \smile \omega_2)(\sigma) = \omega_1([v_0, v_1]) \cdot \omega_2([v_1, v_2]) \pmod 2.
\]
\end{definition}

In a boolean circuit, non-linear gates (AND, NAND, OR) introduce non-linear cochain couplings:
\[
d(\omega_{g_i}) = \delta^0(\omega_{g_i}) \oplus \omega_{g_j} \oplus \omega_{g_k} \oplus (\omega_{g_j} \smile \omega_{g_k}).
\]

\subsection{Formal Statement of Target III}

\begin{target}[Cup-Product Incompressibility]\label{target:3}
Let $K(\phi)$ be the 2-cell complex of a random 3-SAT formula at density $\alpha \in (\alpha_d, \alpha_c)$.

\textbf{Prove that:}
The cup product pairing:
\[
\smile: H^1(K(\phi), \mathbb{F}_2) \times H^1(K(\phi), \mathbb{F}_2) \longrightarrow H^2(K(\phi), \mathbb{F}_2)
\]
has non-trivial rank:
\[
\Rank(\smile) \ge 2^{\Omega(n)}.
\]
Specifically, prove that there is no sequence of linear substitutions (affine coordinate transformations over $\mathbb{F}_2$) of size $\mathrm{poly}(n)$ that maps all non-zero cup products to coboundaries in $\Img(\delta^1)$.
\end{target}

\begin{remark}[Proof Strategy for Target III]
Target III establishes why 3-SAT is fundamentally harder than 2-SAT and Tseitin formulas. 
\begin{itemize}
    \item For \textbf{2-SAT} and \textbf{Tseitin}, the clause constraints are purely linear over $\mathbb{F}_2$, meaning $\omega_j \smile \omega_k = 0$ in cohomology, allowing Gaussian elimination ($\mathbf{P}$-time algorithms) to collapse the complex.
    \item For \textbf{3-SAT}, the non-linear coupling forces $\omega_j \smile \omega_k \neq 0$, creating higher-dimensional topological obstructions that cannot be eliminated by linear matrix inversions.
\end{itemize}
\end{remark}

---

\section{The Unconditional Convergence Theorem}

\begin{theorem}[Conditional Implication to Full Separation]\label{thm:final_convergence}
Suppose Targets \ref{target:1}, \ref{target:2}, and \ref{target:3} hold simultaneously. Then:
\[
\mathbf{NP} \not\subseteq \mathbf{P/poly} \implies \mathbf{P} \neq \mathbf{NP}.
\]
\end{theorem}

\begin{proof}
By Target \ref{target:1}, $\dim H^1(K(\phi), \mathbb{F}_2) \ge 2^{\kappa n}$.
By Target \ref{target:2}, any circuit $C$ deciding 3-SAT must satisfy $\dim \Img(f_C^*) \ge 2^{\kappa n}$.
By the Rank-1 Coboundary Invariant (proven in the Master Monograph, Theorem 7.3), each gate in $C$ can contribute at most $1$ dimension to $\Img(f_C^*)$.
Target \ref{target:3} guarantees that no sub-exponential linear simplification can cancel these dimensions via field transformations.
Therefore:
\[
|C| \ge \dim \Img(f_C^*) \ge 2^{\kappa n} = 2^{\Omega(n)}.
\]
Since this holds for any non-uniform boolean circuit family deciding 3-SAT, 3-SAT $\notin \mathbf{P/poly}$, establishing $\mathbf{NP} \not\subseteq \mathbf{P/poly}$ and consequently $\mathbf{P} \neq \mathbf{NP}$.
\end{proof}

---

\section{Summary: The Exact State of the Proof}

\begin{center}
\begin{tabular}{lll}
\toprule
\textbf{Subsystem / Component} & \textbf{Status} & \textbf{Formal Location} \\
\midrule
Unconditional $\NEXP \not\subseteq \TC^0$ & \textbf{100\% Proven} & Master Monograph, Section 5 \\
Unconditional $\NEXP \not\subseteq \SparseTC$ & \textbf{100\% Proven} & Master Monograph, Section 5 \\
Exact MPS Tensor CAPP & \textbf{100\% Proven} & Master Monograph, Section 3 \\
Rational Chebyshev on DAGs & \textbf{100\% Proven} & Master Monograph, Section 4 \\
Rank-1 Coboundary Update Lemma & \textbf{100\% Proven} & Master Monograph, Theorem 7.3 \\
Zig 0.16.0 SIMD Silicon Suite & \textbf{100\% Verified} & \texttt{src/test\_all\_invariants.zig} \\
\midrule
Target I: Exponential $H^1$ Dimension & \textbf{Open Lemma} & This Document, Section 2 \\
Target II: Functional Surjectivity & \textbf{Open Lemma} & This Document, Section 3 \\
Target III: Cup-Product Non-Triviality & \textbf{Open Lemma} & This Document, Section 4 \\
\bottomrule
\end{tabular}
\end{center}

\end{document}
